\section{Estrategia de despliegue y gobernanza}
\subsection{Estrategia de salida del modelo}
Weston enfatiza que la salida de self-rewarding no debe exponer directamente la raw verified response: es necesario pasar primero por una capa de alignment (meta judge + human checker) que la confirme antes de producir la salida. La estrategia de despliegue suele incluir:
\begin{itemize}
    \item Deployment filters: solo puede ingresar a la downstream API la verified response que pase la verificación de consistencia con el judge.
    \item Conservative fallback: cuando el judge confidence esté por debajo del umbral, se degrada automáticamente al modelo baseline RLHF.
    \item Update schedule: se genera una versión self-rewarding a diario o semanalmente, y se prueba retrospectivamente la hallucination con el offline dataset.
\end{itemize}

\begin{warningbox}{Advertencia de riesgo de despliegue}
Si el modelo self-rewarding presenta hallucination de manera persistente en producción, aunque la verified response haya sido aceptada por el judge, de todos modos puede producir misleading explanations en la conversación. Se necesita un multi-agent cross-check o una nightly calibration para evitar el reward model drift.
\end{warningbox}

\subsection{Gobernanza y auditoría}
A nivel de gobernanza se debe registrar: 1) el judge log de cada verified response; 2) el reward score y el compute cost; 3) la versión del prompt template y de la Self-Instruct seed. Con ayuda de estos log, se puede aportar transparencia para future audits y sustentar la pregunta `Why did the model take this reasoning path?`

\begin{table}[H]
\centering
\begin{tabular}{p{5cm}p{10cm}}
\toprule
Dimensión de gobernanza & Práctica \\
\midrule
Audit trail & Guardar la verified response, los resultados del judge y el reward score para reproducirlos posteriormente. \\
Prompt provenance & Marcar cada Self-Instruct prompt con un seed ID y un generation timestamp. \\
Compute budget governance & Establecer un máximo de tokens para el cross-check; si se excede, se dispara una alert. \\
\bottomrule
\end{tabular}
\caption{La capa de gobernanza hace que el self-improvement sea auditable.}
\end{table}

\subsection{Resumen de la sección}
Las políticas de filtrado previas al despliegue y los registros de gobernanza aseguran que la verified response no genere confusión directamente, y permiten que el equipo pueda responder en cualquier momento ``How did this reasoning chain earn its reward''.

